% ============================================================
% PAPER V3 - SIMON: HYBRID NEURAL FORCE CORRECTION
% Hybrid Neural Force Correction for the Unrestricted Three-Body Problem
% Aarush Gupta
% ST Yau High School Science Award (Asia)
% ============================================================

\documentclass[12pt,a4paper]{article}

% Core packages
\usepackage[a4paper, left=2.5cm, right=2.5cm, top=2.5cm, bottom=2.5cm]{geometry}
\usepackage{graphicx}
\usepackage{amsmath}
\usepackage{amssymb}
\usepackage{float}
\usepackage{setspace}
\usepackage{booktabs}
\usepackage{array}
\usepackage{multirow}
\usepackage{pdfpages}
\usepackage{natbib}
\usepackage{xcolor}
\usepackage{parskip}
\usepackage{caption}
\usepackage[T1]{fontenc}
\usepackage{microtype}

% Hyperref for clickable TOC
\usepackage[
  colorlinks = true,
  linkcolor  = blue,
  citecolor  = blue,
  urlcolor   = blue,
  pdftitle   = {Hybrid Neural Force Correction for the Unrestricted Three-Body Problem},
  pdfauthor  = {Aarush Gupta}
]{hyperref}

\onehalfspacing
\captionsetup{font=small, labelfont=bf, margin=1cm}

\newcommand{\simon}{\texttt{SIMON}}
\newcommand{\ias}{\texttt{ias15}}
\newcommand{\rebound}{\texttt{Rebound}}

\begin{document}

% ============================================================
% COVER PAGE
% ============================================================
\begin{titlepage}
\centering
\vspace*{1.5cm}
{\large \textbf{2026 S.T. Yau High School Science Award (Asia)}}\\[0.5cm]
{\large Research Report}\\[1.5cm]
\hrule
\vspace{0.5cm}
{\small
\begin{tabular}{ll}
\textbf{Registration Number:} & [FILL IN REGISTRATION NUMBER] \\[0.3cm]
\textbf{Name of team member:} & Aarush Gupta \\
\textbf{School:}              & [FILL IN SCHOOL NAME] \\
\textbf{City, Country:}       & [FILL IN CITY, COUNTRY] \\[0.6cm]
\textbf{Name of supervising teacher:} & [FILL IN SUPERVISOR NAME] \\
\textbf{Job Title:}           & [FILL IN JOB TITLE] \\
\textbf{School/Institution:}  & [FILL IN SCHOOL/INSTITUTION] \\
\textbf{City, Country:}       & [FILL IN CITY, COUNTRY] \\
\end{tabular}
}
\vspace{0.5cm}
\hrule
\vspace{2.5cm}
{\large \textbf{Title of Research Report}}\\[1cm]
{\LARGE \textbf{Hybrid Neural Force Correction for the Unrestricted Three-Body Problem}}\\[2.5cm]
{\large \textbf{Date}}\\[0.5cm]
{\large 31 August 2026}
\end{titlepage}

% ============================================================
% TITLE PAGE (Roman page i)
% ============================================================
\pagenumbering{roman}
\setcounter{page}{1}

\begin{center}
{\LARGE \textbf{Hybrid Neural Force Correction for the Unrestricted Three-Body Problem}}\\[0.8cm]
{\large Aarush Gupta}\\[0.3cm]
{\normalsize May 2026}\\[1.2cm]
\end{center}

\begin{abstract}
We present \simon{}, a hybrid neural force correction integrator for the unrestricted three-body gravitational problem, augmenting second-order leapfrog integration with a physically constrained pairwise force magnitude network. Through a systematic four-run ablation study, we demonstrate that \simon{}'s long-horizon orbital stability rests on two orthogonal physical constraints: adaptive sub-stepping, which resolves close-encounter instabilities inherent to fixed-timestep leapfrog integration, and analytic force direction, which prevents the accumulation of non-radial gravitational impulses during close encounters. Removing adaptive sub-stepping while retaining the neural correction causes body ejection to $161.7$\,AU; replacing analytic direction with a learned force direction causes ejection to $83.0$\,AU despite adaptive stepping remaining active. The neural network correction itself is shown to be an accuracy mechanism rather than a stability mechanism: removing it while retaining both physical constraints maintains bounded orbital dynamics with an 8.5\% larger final position error ($1.67$ versus $1.54$\,AU over 100 years for the default initial condition). We evaluate \simon{} across seven initial conditions spanning the full dynamical spectrum from near-circular to strongly-interacting configurations, confirming bounded dynamics for four physically stable configurations and correct reproduction of instability for three physically unstable ones. Energy drift is attributed to the second-order leapfrog integrator rather than the neural correction, confirmed quantitatively by timestep scaling with fitted slopes of $1.32$--$2.43$ across configurations (theoretical prediction: $2.0$). Training reproducibility is confirmed across five random initialisations, with the divergence-rate proxy $\lambda = 0.1655\,\mathrm{yr}^{-1}$ invariant across all seeds. At the operational timestep $dt = 0.04$\,yr, \simon{} achieves a $1.10{\times}$ speedup over the high-precision \ias{} integrator while maintaining bounded three-body dynamics over 100-year horizons.
\end{abstract}

\bigskip
\noindent\textbf{Keywords:} Three-Body Problem, Hybrid Neural Force Correction, Leapfrog Integration, Adaptive Sub-Stepping, Lyapunov Divergence, Orbital Stability, Ablation Study, Chaotic Dynamics

\newpage

% ============================================================
% ACKNOWLEDGEMENT (Page ii)
% ============================================================
\section*{Acknowledgement}
\addcontentsline{toc}{section}{Acknowledgement}

\noindent [This section should acknowledge all assistance and contributions received during the research, including the supervisor, institutions, funding sources, and any computational resources used. Please fill in before submission.]

\newpage

% ============================================================
% ACADEMIC HONESTY (Page iii)
% Upload Declaration_of_academic_integrity.pdf to Overleaf
% ============================================================
\includepdf[pages=-]{Declaration_of_academic_integrity.pdf}

% ============================================================
% TABLE OF CONTENTS (Pages iv-v)
% ============================================================
\tableofcontents
\newpage

% ============================================================
% MAIN REPORT -- Arabic page numbering
% ============================================================
\pagenumbering{arabic}
\setcounter{page}{1}

% ============================================================
% SECTION 1 -- INTRODUCTION (NEW)
% ============================================================
\section{Introduction}
\label{sec:intro}

The Three-Body Problem---predicting the long-term gravitational evolution of three mutually interacting masses---has occupied celestial mechanics for over three centuries. Unlike the two-body case, which admits closed-form analytical solutions, the three-body system is generically chaotic: trajectories are exponentially sensitive to initial conditions, rendering exact long-horizon prediction impossible. Despite this, the qualitative dynamical structure of such systems---whether orbits remain bounded, when close encounters occur, and at what rate trajectories diverge---carries significant scientific value for planetary dynamics, stellar cluster modelling, and the stability of exoplanetary systems.

Traditional high-precision integrators, such as the 15th-order adaptive Gauss-Radau method \ias{} \citep{Rein_2014}, achieve near-machine-precision energy conservation but scale computationally with simulation duration. Neural network-based approaches offer faster force evaluations; however, standard networks fail to enforce the physical constraints necessary for long-horizon stability, particularly during close encounters where the gravitational force diverges as $r \to 0$.

This paper presents \simon{}, a hybrid integrator that augments second-order leapfrog integration with a physically constrained pairwise force magnitude network. In V2 of this work, \simon{} was shown to achieve bounded long-horizon dynamics at computational speedups of up to $2.07\times$ over \ias{}. That work, however, left open a fundamental question: which components of \simon{}'s architecture are actually responsible for orbital stability?

This V3 paper addresses that question directly through a systematic four-run ablation study, yielding the central finding of this work: \textbf{\simon{}'s orbital stability rests on two orthogonal physical constraints}. Adaptive sub-stepping resolves the close-encounter instabilities that would otherwise cause body ejection under fixed-timestep leapfrog integration. Analytic force direction prevents the accumulation of non-radial impulses that cause ejection even when adaptive stepping is active. The neural network correction, by contrast, is an accuracy mechanism rather than a stability mechanism: removing it maintains bounded dynamics with only an 8.5\% larger final position error.

The contributions of this paper are as follows:
\begin{itemize}
    \item A four-run ablation framework decomposing \simon{}'s stability into two orthogonal mechanisms (Section~\ref{sec:results_stab}).
    \item Evaluation across seven initial conditions spanning the full dynamical spectrum, with honest framing of physically unstable cases confirmed independently by \ias{} (Section~\ref{sec:results_gen}).
    \item Energy conservation analysis attributing drift to the leapfrog integrator rather than the neural correction, confirmed by $dt^2$ scaling (Section~\ref{sec:energy}).
    \item Training reproducibility across five random initialisations (Section~\ref{sec:seeds}).
\end{itemize}

The paper is structured as follows. Section~\ref{sec:litreview} reviews the Three-Body Problem and relevant literature. Section~\ref{sec:comp_modeling} discusses computational modelling. Section~\ref{sec:methodology} describes the \simon{} architecture. Section~\ref{sec:results} presents experimental results. Section~\ref{sec:discussion} provides discussion and conclusions.

% ============================================================
% SECTION 2 -- LITERATURE REVIEW (V2 Section 1 VERBATIM)
% ============================================================
\section{Literature Review: The Three-Body Problem and Neural Approaches}
\label{sec:litreview}

The Three-Body Problem is the classical problem in celestial mechanics of describing the motion of three massive bodies gravitationally interacting with each other. This problem has significant applications in understanding the orbits of planets, moons, and artificial satellites \citep{gardner2006james}.

The Three-Body Problem is the manifestation of centuries of development in astronomy. Its history dates back to ancient Greek times. The Greeks were among the first to begin the systematic observation of celestial movements, although they lacked the mathematical framework to describe the complex gravitational interactions between bodies precisely \citep{anjum2020timelinegravity}. However, in the 17th century, Sir Isaac Newton's \textit{Philosophiæ Naturalis Principia Mathematica} \citep{newton1687principia, cohen1999principia} dictated universal laws that laid the groundwork for classical mechanics, creating a wide range of avenues for research into gravitational systems. To utilize Newton's formalism to understand the Three-Body Problem, we first need to analyze the simpler Two-Body Problem.

\subsection{The Two-Body Problem}
\label{subsec:2bp}
The general formula for the force of attraction between any two objects in a gravitational field is given by Newton's Universal Law of Gravitation in Equation~\ref{eq:1}:

\begin{equation}
    F=G\frac{m_1m_2}{r^2},
    \label{eq:1}
\end{equation}
where $G$ is the universal gravitational constant, $m_1$ is the mass of the first body, $m_2$ is the mass of the second body and $r$ is the distance between the centers of the bodies. Using this formula, the total force of gravity experienced by an object from any number of massive bodies can be calculated.

To understand how this force translates to celestial orbital motion, consider the following scenario: an observer on earth throws a ball of mass $m_b$. The harder the ball is thrown, the farther it will travel, but will eventually fall to the ground as a result of gravitational attraction exerted on it by the earth. However, if thrown hard enough, the ball will make an orbit around the earth and return to where it started. If we assume that the ball moves in a circular orbit around the Earth, we can equate the force of gravity exerted on the ball by the earth to the centripetal force on the ball, giving us Equations~\ref{eq:2}, and~\ref{eq:3}.

\begin{equation}
    G\frac{m_e m_b}{r^2} = \frac{m_b v^2}{r}
    \label{eq:2}
\end{equation}

\begin{equation}
    v= \sqrt{\frac{G m_e}{r}},
    \label{eq:3}
\end{equation}

where $m_e$ is the mass of the Earth, $v$ is the velocity of the ball, $G$ is the universal gravitational constant, and $r$ is the distance between the two bodies, measured from their centers of mass.

\begin{figure}[H]
    \centering
    \includegraphics[width=0.5\linewidth]{Figure_1.pdf}
    \caption{A diagram of a simple two-body system.}
    \label{fig:fig1}
\end{figure}

Similarly, a two-body system (see Figure~\ref{fig:fig1}) can be reduced to operate under the same principles as the ball and the earth example. Therefore, the velocity of the orbiting planet can be given by Equation~\ref{eq:3}. In the math presented above, the path of the ball around the Earth is approximated to be circular. However, in reality, the ball can take any elliptical path around the Earth \citep{murray1999solar}. The two-body framework can be used to model systems such as a satellite orbiting the earth or a planet orbiting the sun. In addition, the two-body problem is helpful in modeling N-body systems where the masses of two objects are much larger than the masses of all other bodies. However, the two-body problem omits situations in which three or more masses are all exhibiting comparable gravitational forces on each other at the same time.

\subsection{The Three-Body Problem}
\label{subsec:ODEs}
In the three body problem, the goal is to approximate the future positions and velocities of the \textit{three} planets, given their initial states. However, when calculating the future states of a system of three bodies interacting gravitationally, it is not possible to construct an equation that will precisely reveal the trajectories of each body. To obtain this conclusion, we define some assumptions and simplifications.

\begin{itemize}
    \item \textbf{Point Masses:} The bodies will be treated as point masses.
    \item \textbf{Newtonian Gravity:} We will use Newton's Law of Gravitation to describe the force of gravity.
    \item \textbf{Isolated System:} The system only consists of the three bodies interacting, without external forces.
\end{itemize}

We present a schematic of this system in Figure~\ref{fig:fig2}.

\begin{figure}[H]
    \centering
    \includegraphics[width=0.5\linewidth]{Figure_2.pdf}
    \caption{A diagram of a simple three-body system. Each body is exerting a gravitational force on the other two bodies.}
    \label{fig:fig2}
\end{figure}

We will use a three-dimensional cartesian coordinate system to define the centers of each of the bodies. Let the position vector of a body $n$ be the vector $\mathbf{P}_n$. Each of the six forces in the system can be described using Newton's law of gravitation, generally given by
$$
\mathbf{F}_{n/b} = G \frac{M_n M_b}{\left|\mathbf{P}_n - \mathbf{P}_b\right|^3} (\mathbf{P}_n - \mathbf{P}_b)
$$
Where $\textbf{F}_{n/b}$ is the gravitational force of attraction that body $n$ feels toward body $b$. This is the first law we will use to formulate our problem. We will also take Newton's second law as follows.
$$\sum \mathbf{F}_{n} = m\mathbf{a_n}$$
$$\sum \mathbf{F}_{n} = m \left( \frac{d\mathbf{v}_n}{dt} \right)$$
$$\sum \mathbf{F}_{n} = m \left( \frac{d^2 \mathbf{P}_n}{dt^2} \right)$$
We can now use these two laws to create a system of Ordinary Differential Equations (ODEs) that we can use to approximate the positions of the bodies at time $t$. We can take the sum of forces on body 1 as an example:
$$\sum \mathbf{F}_{\text{exterior}} (\text{body } 1) = m_1 \frac{d^2\mathbf{P}_1}{dt^2} = \mathbf{F}_{2/1} + \mathbf{F}_{3/1}$$
Substituting in law 1:
$$\sum \mathbf{F}_{\text{exterior}} (\text{body } 1) = G \left[ M_2 \frac{\mathbf{P}_2 - \mathbf{P}_1}{|\mathbf{P}_2 - \mathbf{P}_1|^3} + M_3 \frac{\mathbf{P}_3 - \mathbf{P}_1}{|\mathbf{P}_3 - \mathbf{P}_1|^3} \right]$$
$$\frac{d^2\mathbf{P}_1}{dt^2} = G \left[ M_2 \frac{\mathbf{P}_2 - \mathbf{P}_1}{|\mathbf{P}_2 - \mathbf{P}_1|^3} + M_3 \frac{\mathbf{P}_3 - \mathbf{P}_1}{|\mathbf{P}_3 - \mathbf{P}_1|^3} \right]$$

Doing the same for bodies 2 and 3 we get the system of ODEs \citep{murray1999solar}:
\begin{align*}
\frac{d^2\mathbf{P}_1}{dt^2} &= G \left[ M_2 \frac{\mathbf{P}_2 - \mathbf{P}_1}{|\mathbf{P}_2 - \mathbf{P}_1|^3} + M_3 \frac{\mathbf{P}_3 - \mathbf{P}_1}{|\mathbf{P}_3 - \mathbf{P}_1|^3} \right] \\[1em]
\frac{d^2\mathbf{P}_2}{dt^2} &= G \left[ M_1 \frac{\mathbf{P}_1 - \mathbf{P}_2}{|\mathbf{P}_1 - \mathbf{P}_2|^3} + M_3 \frac{\mathbf{P}_3 - \mathbf{P}_2}{|\mathbf{P}_3 - \mathbf{P}_2|^3} \right] \\[1em]
\frac{d^2\mathbf{P}_3}{dt^2} &= G \left[ M_1 \frac{\mathbf{P}_1 - \mathbf{P}_3}{|\mathbf{P}_1 - \mathbf{P}_3|^3} + M_2 \frac{\mathbf{P}_2 - \mathbf{P}_3}{|\mathbf{P}_2 - \mathbf{P}_3|^3} \right]
\end{align*}

This system of ODEs does not have a closed-form analytical solution like the two-body example in Section~\ref{subsec:2bp}. This is because the motion of the three bodies that all exert forces simultaneously on each other is too complex and generally leads to chaos. In this context, chaos means that even the slightest change in position or velocity has a relatively huge impact on the long-term motion of the system, leading to unpredictable motion which makes it impossible to measure with complete accuracy.

\subsection{Chaos}
\label{Chaos}
Chaos can be quantified using many different frameworks. For the sake of simplicity, the framework we will be using is the Lyapunov Exponent \citep{ott2002chaos} which measures the rate of trajectory divergence as an exponential model. If the initial separation between two points in phase space is $\delta(0)$, their separation at time $t$ is approximated by the function
$$\delta(t) \approx e^{\lambda t} \delta(0)$$
where $\lambda$ is the Lyapunov exponent. If $\lambda > 0$, the system is chaotic, and if $\lambda \leq 0$ it is regular, or periodic. $\lambda$ is a time constant.

The ideas of chaos in the Three-Body Problem can be traced back to the late 19th century, when Henri Poincar\'{e} introduced the groundwork for defining Chaos \citep{poincare1890sur}. He found that when perturbations (deviance due to an external force) in a chaotic system occurred, the stable and unstable manifolds (two different sets of future states of a system, one stable versus unstable) of a periodic orbit began to intersect. Instead of one manifold neatly returning to its source, they crossed each other an infinite number of times. Because they cannot cross themselves but must cross each other to satisfy the equations of motion, they create an incredibly intricate, mesh-like structure. This tangling, known as the homoclinic tangle, leads to chaos because it creates an infinite, fractal-like web of trajectories, that makes the system's future state impossible to analytically predict.

\begin{figure}[H]
    \centering
    \includegraphics[width=0.5\linewidth]{homoclinic.pdf}
    \caption{A plot of a homoclinic tangle.}
    \label{fig:fig3}
\end{figure}

Poincar\'{e} himself noted he couldn't even begin to draw the figure, concluding that an algebraic formula to predict a planet's long-term position was physically impossible. Linking back to the Lyapunov exponent, Poincar\'{e} predicted that the rate of trajectory divergence would be exponential. Two bodies starting at almost identical positions within the tangle will be stretched along the unstable manifold and compressed along the stable one. Very quickly, their paths smear exponentially across the web.

This work showcased that the Three-Body problem is highly sensitive to its initial conditions, exemplifying how seemingly simple gravitational systems can exhibit unpredictable behavior. This would go on to establish the idea that while some solutions existed, a generalized solution to the Three-Body Problem was impossible, which set up the connection between chaos theory and celestial mechanics.

Some specific factors that influence the chaos of a Three-Body system are mass, relative velocity, and initial separation. For example, if the mass of a system is concentrated on one body while the other two are smaller, then the system is less likely to be as chaotic than if the masses were distributed \cite{Hietarinta_1993}. Furthermore, if the relative velocities of the bodies are high, even when nearby, the bodies may be moving with too much kinetic energy to induce an exchange that would cause them to diverge much from their initial trajectories. This would reduce the rate at which changes in trajectory would happen, lowering $\lambda$. Lastly, if the bodies have a large initial separation, it is unlikely that they will enter each other's zones of gravitational influence, so the likelihood of divergence from their trajectories will be lower, reducing $\lambda$ and hence the chaos of the system.

\subsection{The Restricted Three-Body Problem}
There exists another form of the Three-Body Problem by the name of the Restricted Three-Body Problem, which was developed following Newton's discoveries primarily by Euler and Lagrange during the 18th century. Euler built upon the first versions of the restricted Three-Body problem by fixing two massive bodies in space and solving analytically for the motion of a third, smaller particle under their combined gravitational pull \cite{bistafa2021euler}. This makes the problem easier to interpret, with one notable method of interpretation being Lagrange points. Developed by Joseph-Louis Lagrange, these are five special solutions which represent stable positions where a smaller body can theoretically maintain its position \cite{lagrange1772essai}. However, these solutions tended to be unstable in the context of the Unrestricted Three-Body problem.

The fact that these stable Lagrange points can be identified means that the restricted Three-Body Problem tends to be more applicable than the more chaotic Unrestricted version, because the restrictions simplify the system enough to be studied to a great degree of accuracy without losing too much of the complex dynamics. For example, the restricted Three-Body Problem is used to chart spacecraft trajectories when entering the gravitational influence of celestial bodies. Moreover, identifying Lagrange points in a restricted Three-Body System allows telescopes like the James Webb telescope \citep{gardner2006james} to stay in an orbit to perform long-term observational missions.

\subsection{Expansion into Modern Physics}
\label{exp_phys}
Progressing into the 20th century, the Three-Body problem continued to attract significant attention. The development of new numerical methods brought new perspectives on chaos and stability in the Three-Body problem. Furthermore, the development of modern computers allowed for numerical simulations of the Three-Body problem via integration, allowing easier investigation into its chaotic nature for longer time periods.

In modern astrophysics, the Three-Body problem finds application in contexts like calculating orbits of planets in multi-body systems, examining interactions in exoplanet systems, and understanding stability in star clusters. On a more recent note, as this paper will outline, machine learning techniques are being used to deepen traditional methodologies of Three-Body analysis and prediction. One such example of a machine learning application was by \cite{ulibarrena2024hybrid}. In their setup, they have created and evaluated a Hamiltonian Neural Network (HNN,) and their findings were that an HNN was able to successfully emulate a traditional step-by-step integrator, at a considerably reduced computational cost. A Hamiltonian Neural Network is a type of NN that composes forces in a way such that energy conservation is structurally preserved. However, this increased complexity was reflected in their results, as the training process was significantly more advanced than that of other methods of Three-Body computation. The Three-Body Problem continues to be an area of interest for modern computational physics problems, and its chaotic nature makes it especially beneficial as a testbed for numerical stability in physics applied machine learning.

\subsection{Summary}

To summarise, we have defined:

\begin{itemize}
    \item \textbf{The Two-Body Problem:} Assuming circular orbits, using Newton's universal law of gravitation.
    \item \textbf{The Three-Body Problem:} The setup and derivation of a system of ODEs for approximating the motion of an unrestricted Three-Body system.
    \item \textbf{Chaotic Nature of a Three-Body System:} Introduced the idea of Chaos using the Lyapunov exponent (rate of trajectory divergence), and established that the trajectories in a Three-Body system are highly sensitive to the initial conditions, rendering it impossible to predict without approximation. Explored Poincar\'{e}'s discoveries on homoclinic tangles and their implications on chaos in the Three-Body Problem.
    \item \textbf{Examples of Influences of Initial Conditions on the Chaos of a Three-Body System:} Explained that a system with more distributed masses, closer bodies, and lower relative velocities is likely to be highly chaotic.
    \item \textbf{The Restricted Three-Body Problem:} Analysis of Lagrange points, and how they are applied to real-life scenarios.
\end{itemize}

% ============================================================
% SECTION 3 -- COMPUTATIONAL MODELING (V2 Section 2 VERBATIM)
% ============================================================
\section{Computational Modeling of the Three-Body Problem}
\label{sec:comp_modeling}
The Gravitational Three-Body Problem, can be simulated using the aid of computer programs, as mentioned in Section~\ref{exp_phys}. Because the system of ODEs defining the Unrestricted Three-Body Problem does not have a closed solution (see Section~\ref{sec:litreview}), integration is done incrementally to calculate the evolution of the system over a set period of time. In short term contexts, traditional integration techniques tend to work quite well and provide results to a great level of accuracy. However, traditional integration methods scale linearly with the integration period. Hence, it can be useful to find a technique that doesn't scale with the integration period for longer term simulations so as to be more computationally efficient. Furthermore, in a larger time frame, errors may accumulate as a byproduct of the highly chaotic nature of a Three-Body system.

A Neural Network may offer a method of simulating a system over longer integration periods more efficiently. By training on datasets with a variety of trajectories, neural networks can learn to directly map from a set of initial conditions to a future state. This allows it to skip the incremental time steps that a traditional integrator would have to take, which allows the network to perform fewer calculations, and would ideally increase the speed of a prediction. Moreover, fewer calculations may mean less floating-point errors, which could be beneficial for stability in long-term simulations. The computing cost that comes with using a Neural Network that maps directly from initial conditions to a future state would remain constant, regardless of the temporal distance $t$ from the starting point, which in principle could yield large speedups over step-by-step integrators. In practice, however, the hybrid approach in this work still evolves the system through timesteps, so realized speedups depend on the specific architecture and timestep chosen. However, a standard network may fail to acknowledge the underlying physics that govern the Three-Body problem, which may lead to energy conservation violations. Furthermore, it may struggle during close encounters. This is because according to Newton's law of universal gravitation, the gravitational force exerted between two celestial bodies separated by distance $r$ is proportional to $1/r^2$. Therefore, when $r$ approaches $0$ during a close encounter, the gravitational force approaches $\infty$. This relation has such a steep gradient that even the slightest error in position from the network may cause a body to be ejected from the system or gain infinite energy, which is not physically possible \cite{Dehnen2001}.

These shortcomings render computational methods of simulating the Three-Body problem incomplete, and for this reason, we will study the use of more complex networks.

% ============================================================
% SECTION 4 -- METHODOLOGY (V2 Section 3 VERBATIM)
% ============================================================
\section{SIMON: Architecture and Methodology}
\label{sec:methodology}

\subsection{Overview and Research Objective}
\label{research_obj}
The aim of this study is to investigate how closely physically constrained neural network force models can reproduce the long-term stability and energy-conserving behavior of traditional numerical integrators in three-body simulations, while operating at significantly reduced computational cost. Rather than attempting to replace established integration schemes, this work seeks to quantify the extent to which neural networks can approximate the physical fidelity of traditional step integrators when computational efficiency is prioritized.

Traditional numerical integrators are expected to achieve superior accuracy, particularly with energy conservation over long timescales, but typically incur higher computational costs. Neural methods, by contrast, offer faster force evaluations but often suffer from less long-term stability. By embedding physical constraints directly into the neural network architecture, this study examines whether such models can narrow the gap in physical accuracy while retaining their computational advantage.

Formally, we address the following research question:
\begin{quote}
\emph{How closely can physically constrained neural network force models approximate the long-term stability of traditional numerical integrators in three-body simulations, while operating at substantially higher computational speed?}
\end{quote}

\subsection{Problem Formulation}
We study the unrestricted Three-Body problem where masses solely interact through Newtonian gravitational forces. The governing equations (see Subsection~\ref{subsec:ODEs}) for the evolution of this system are deterministic, time-reversible, and conserve total energy. As discussed earlier in section Subsection~\ref{subsec:ODEs}, a Three-Body system exhibits chaotic motion for generic initial conditions, making it ideal for testing long-term numerical stability in Physically constrained neural networks.

\subsection{Baseline Numerical Solver: \texttt{Rebound}}
Throughout this paper, we will take the ground truth of the Three-Body integration to be the output from \texttt{Rebound} \citep{Rein2012}, and any output from neural networks will be compared against \texttt{Rebound} output. \texttt{Rebound} is an N-body integrator that can simulate the motion of particles under the influence of gravity. \texttt{Rebound} has several integrators: and the one we will be using in our case will be the \texttt{ias15}, \cite{Rein_2014} as it is one of the most precise integrators offered by Rebound. It is a high-order adaptive integrator based on Gauss-Radau quadrature, accurate to near machine precision over long integrations \citep{Rein_2014}. The ``15'' in \texttt{ias15} refers to its 15th-order accuracy, which enables excellent energy conservation without being a symplectic method. Its adaptive timestep control makes it especially reliable for close encounters and long-horizon simulations.

To get an idea of what the model we will train is actually being trained on, observe 2 simulations with trajectories plotted.

\begin{figure}[H]
    \centering
    \includegraphics[width=0.7\linewidth]{rebound_sim1.pdf}
    \caption{Trajectories calculated by \texttt{Rebound}, with body masses $1\,M_\odot$, $0.01\,M_\odot$, $0.005\,M_\odot$, for 100 years.}
    \label{fig:rebound1}
\end{figure}

\begin{figure}[H]
    \centering
    \includegraphics[width=0.7\linewidth]{rebound_sim2.pdf}
    \caption{Trajectories calculated by \texttt{Rebound}, with body masses $1\,M_\odot$, $0.01\,M_\odot$, $0.005\,M_\odot$, for 100 years, but with different initial positions.}
    \label{fig:rebound2}
\end{figure}

These simulations serve several purposes. First, they establish a physically consistent baseline against which alternative force models can be evaluated. Second, they illustrate the range of orbital behaviors encountered even in relatively simple three-body configurations, including non-circular motion and mutual perturbations. Finally, it provides a concrete visualization of the types of interactions, particularly variations in inter-body separation and acceleration, that any force emulator must accurately capture in order to remain stable over extended integration times.

% V3 NOTE: New paragraph added -- ias15 adaptive stepping context
\noindent\textit{[V3 addition: note on \ias{} adaptive stepping]} \ias{} uses its own built-in continuous global adaptive timestepping at every integration step, adjusting its internal step size based on a local error estimate. This is why \ias{} achieves near machine-precision energy conservation regardless of the nominal $dt$ parameter. \simon{}'s adaptive sub-stepping (Section~\ref{subsec:adaptive}) is fundamentally different: it is targeted and local, activating only when a close encounter is detected ($r_{\min} < 0.05$\,AU), leaving all other steps at the base fixed timestep.

\subsection{Physically Constrained Force Magnification Neural Network}
\label{subsec:pfmn}
The Pairwise Force Magnitude Network (PFMN) \cite{nguyen2021pairwise}, \cite{battaglia2016interaction}, we are training for the purpose of this investigation is not trained to predict forces directly, but instead is trained to predict a scalar pairwise force magnitude, which is essentially the magnitude of the force of interaction between bodies interacting, given in pairs. Let:
\begin{itemize}
    \item \textbf{$\mathbf{x}_i(t) \in \mathbb{R}^3$} be the position vector of body $i$ at a time $t$,
    \item \textbf{$\mathbf{v}_i (t) = \mathbf{\dot x}_i(t)$} be its velocity vector,
    \item \textbf{$m_i > 0$} be its mass,
    \item \textbf{$G$} be the universal gravitational constant.
\end{itemize}

For any pair of bodies, the displacement vector can be given as $r_{ij} = x_j - x_i$. The pairwise distance can be given by $r_{ij} = ||r_{ij}||$. The classical Newtonian equations of motion used for their interactions are; $\sum \mathbf{F}_{ij} = m_i \mathbf{\ddot x}$ (second law of motion,) and $$\mathbf{F}_{ij}^{\text{Newt}} = G \frac{m_i m_j}{|r_{ij}|^3} \mathbf{r}_{ij}$$ (universal law of gravitation.) \cite{murray1999solar}.

To improve stability during close encounters, a softened gravitational force is also computed using a fixed softening parameter $\epsilon$. This softened force replaces the Newtonian inter-body distance $r_{ij}$ with $\sqrt{r^2_{ij} + \epsilon^2}$ to ensure forces are controlled during interactions where the inter-body distance $r_{ij}$ is less than $\epsilon$. \cite{Dehnen2001}. This provides a smoother learning target for the model, in the sense that minuscule errors in position do not completely break the simulation. The new, softened force can be represented as $\tilde{r}$. However, when using softening, a slight bias is introduced, which why during model analysis/evaluation we will be using unsoftened forces.

Hence, substituting into the universal law of gravitation, we get the following representation of a pairwise force:
$$\mathbf{F}_{ij}^{\text{soft}} = G \frac{m_i m_j}{\tilde{r}_{ij}^3} \mathbf{r}_{ij}.$$

Now that the mathematics of the method of force formulation has been developed, we will now move onto the model architecture. The learned model is a scalar function with neural network weights $\theta$, $f_{\theta} : (\tilde{r}_{ij}, m_i, m_j) \mapsto F_{ij} \in \mathbb{R}$, where $\hat{F}_{ij}$ is the predicted pairwise force magnitude. When we implement this, the network will receive a feature vector $z_{ij} = [r_{ij}, m_i, m_j] \in \mathbb{R}^3$, and will output $\hat{F}_{ij} = f_{\theta}(z_{ij})$. We now construct a unit vector for the direction of the force, given by $u_{ij} \equiv \frac{\mathbf{r}_{ij}}{\|\tilde{r_{ij}}\|}$. Now that we have both the direction of the unit vector and the magnitude, we can multiply $\hat{F}_{ij}$ and $u_{ij}$ to get the magnitude and direction of the force: $\mathbf{F}_{ij}^{NN} = F_{ij} \mathbf{u}_{ij}$. Now, to get the sum of the forces on body $i$ we will sum over all forces: $\mathbf{F}_{i}^{NN} = \sum_{j \neq i} \mathbf{F}_{ij}^{NN}$. \cite{kipf2018neural}.

Moreover, the architecture is a neural network with 3 hidden layers of 32 neurons each, using SiLU activation functions:
$\text{Input}(3) \to \text{Linear}(32) \to \text{SiLU} \to \text{Linear}(32) \to \text{SiLU} \to \text{Linear}(32) \to \text{SiLU} \to \text{Linear}(1)$.

% V3 NOTE: New paragraph added -- force direction clarification
\noindent\textit{[V3 addition: force direction]} In the implementation, the force direction unit vector is computed using the unsoftened separation: $u_{ij} = \mathbf{r}_{ij} / \|\mathbf{r}_{ij}\|$. Softening appears only in the force magnitude denominator $\tilde{r}_{ij}^3$. This analytic, exact force direction is a critical physical constraint: Section~\ref{subsec:run_b} demonstrates that replacing it with a learned direction causes body ejection even when all other components of \simon{} are active.

\subsection{Hybrid Incorporation}
\cite{ulibarrena2024hybrid} predicts that using a hybrid combination of an analytical integrator (like \texttt{ias15} in our case) to correct inaccurate neural network predictions could ``prevent large energy errors.'' A fixed timestep is used. We incorporate \texttt{ias15} into our model in a similar way \cite{ulibarrena2024hybrid} did, replacing inaccurate neural predictions with an analytic one instead. We do this by not letting the NN predict by itself, but rather learn a correction, $c$, to a softened gravitational force \cite{kipf2018neural}. The coefficient of correction, predicted by the NN, determines whether the softened force should be stronger or weaker. We divide force computation into three tiers. Far field, where $r > 500\epsilon = 0.15\,\text{AU}$ means the NN is not invoked and is identical to what \texttt{ias15} predicts. Close encounters, where $r < 0.15\,\text{AU}$ means the NN is invoked and predicts a correction factor. To prevent the force from being strengthened or weakened too much, we gate the coefficient, bounding it between 0.2 and 5. If it falls outside of this range, we do not use the NN correction and simply take the softened baseline force. Lastly, for very close encounters, where $r < 5\times10^{-4}$ the NN is again bypassed to ensure that no bodies are ejected, and non-physical behavior is prevented.

% V3 NOTE: New paragraph added -- softening justification
\noindent\textit{[V3 addition: softening consistency]} At the distances where the NN correction is operationally applied ($r \geq 5\times10^{-4}$\,AU), the difference between softened and unsoftened force magnitudes is less than 0.01\%, making the training-evaluation inconsistency negligible in practice.

\subsection{Training Process}
\label{subsec:training}
Training data is generated by sampling 50,000 physically valid pairwise configurations from \texttt{Rebound}. For each configuration, reference forces are computed analytically using Newtonian gravity. Approximately $60\%$ of the samples are drawn from regions where the inter-body distance $r\in [0.5 \epsilon, 50 \epsilon]$ ($\epsilon = 3 \times10^{-4}$), where the correction predicted by the model would deviate meaningfully from 1. The other $40\%$ is sampled from the far field $r\in[50\epsilon, 10]$ where the correction would be closer to one.

Furthermore, the following hyperparameters are used for the model.
\begin{itemize}
    \item \textbf{Optimizer}: AdamW, weight decay = $1\times10^{-4}$
    \item \textbf{Learning rate}: $1\times10^{-3}$, cosine annealing to $1\times10^{-5}$ \cite{loshchilov2017sgdr} over 5,000 epochs.
    \item \textbf{Loss Function}: MSE on $\log(c)$
    \item \textbf{Batch Size}: Full batch (40,000 training samples)
    \item \textbf{Train/evaluation split}: 80\% / 20\% (40,000/10,000)
    \item \textbf{Training device}: NVIDIA GeForce RTX 5050 CUDA 13.0
    \item \textbf{Training time}: 9 seconds
    \item \textbf{Best MSE}: $1.1 \times 10^{-5}$
\end{itemize}

\begin{figure}[H]
    \centering
    \includegraphics[width=0.75\linewidth]{training_log.png}
    \caption{Training and validation loss curves over 5,000 epochs for the \simon{} PairCorrectionNN model (seed 42). Validation MSE converges to $1.1\times10^{-5}$. [V3 addition]}
    \label{fig:training_log}
\end{figure}

\subsection{Adaptive Sub-Stepping}
\label{subsec:adaptive}
The integrator uses a fixed macro timestep $dt$. To resolve close encounters where gravitational acceleration changes faster than $dt$ can capture, an adaptive sub-stepping mechanism is applied \cite{Rein_2014}. At each step $dt$ the minimum pairwise distance $r_{min}$ is computed across all body pairs. If $r_{min} <0.05\,\text{AU}$ then $dt$ will be split into smaller sub-steps. The number of sub-steps is determined by the following: $n_{sub}=\frac{0.05}{r_{min}}$. Each sub-step has a duration $\frac{dt}{n_{sub}}$. This makes sure that as encounters become closer and closer, the number of sub-steps increases proportionally to account for the increased gravitational force (and hence acceleration) that is associated with closer encounters. The maximum of 16 sub-steps per large $dt$ caps the computational cost to optimize speed.

All physics (positions, velocities, accelerations, forces) is computed in float64 for numerical stability. The NN forward pass uses float32 weights, with the correction factor being converted back to float64 before use.

\subsection{Evaluation Strategy}
\label{subsec:eval}
Long-term trajectory divergence in chaotic systems like the Three-Body problem is inevitable, as established in section Subsection~\ref{subsec:ODEs}. With this knowledge in mind, we will evaluate with the following diagnostics on the hybrid model:

\begin{itemize}
    \item \textbf{Trajectory Overlay:} Trajectories generated by neural models are superimposed against the reference numerical solution to gain qualitative insight into the accuracy of the model. This will be done with various different initial conditions to explore how the model will handle different situations, to get a better insight into its versatility. Moreover, we will also plot a time series overlay to analyze if there is any phase drift.
    \item \textbf{Trajectory Divergence:} To gain insight into the sensitivity of each system with respect to the initial conditions, we numerically estimate the trajectory divergence by evolving predicted trajectories in parallel, then calculating the RMS position error between the two. Agreement between neural and reference values indicates preservation of the system's chaotic structure.
    \item \textbf{Speed-Accuracy Tradeoff:} The distance between the neural network predicted trajectory and the \texttt{Rebound} predicted trajectory, plotted over throughput (how fast trajectory points were saved.)
\end{itemize}

% V3 NOTE: New paragraph added -- lambda disclaimer
\noindent\textit{[V3 addition: divergence-rate proxy]} The fitted slope $\lambda$ from the 10--50 year window is a divergence-rate proxy rather than a formal Lyapunov exponent: we compare two different models (\simon{} and \ias{}) rather than two nearby initial conditions evolved under the same dynamics. Consequently, $\lambda$ reflects both the intrinsic chaotic sensitivity of the three-body system and \simon{}'s accumulated force errors. A perfect model would yield $\lambda = 0$. All $\lambda$ values reported in this paper are computed at the operational timestep $dt = 0.04$\,yr using log-linear regression over the 10--50 year window, giving $\lambda = 0.1655$\,yr$^{-1}$ for the default initial condition (IC1).

\subsection{Scope of the Method}
This method is designed to see how closely a physically constrained Neural Network can shadow a traditional iterative integrator in terms of close encounters, computational efficiency, and trajectory divergence. By comparing these two, we seek insight into the development of neural methods of Three-Body emulation, and learn to what extent these methods are able to speed-up simulations without compromising significant accuracy and stability.

% ============================================================
% SECTION 5 -- RESULTS (NEW)
% ============================================================
\section{Results}
\label{sec:results}

In this section, we examine the results of the comparison between \texttt{Rebound} and \simon{}. We will be seeking to understand the results through a numerical and intuitive perspective, and use the insights gained to answer the research question in Section~\ref{research_obj}. Results are organised into three parts: stability decomposition via ablation (Section~\ref{sec:results_stab}), generalisation across initial conditions (Section~\ref{sec:results_gen}), and physical consistency with reproducibility (Section~\ref{sec:results_phys}).

% ---- 5.1 RESULTS I: STABILITY DECOMPOSITION ----
\subsection{Results I: Stability Decomposition}
\label{sec:results_stab}

\subsubsection{The Four-Run Ablation Design}
\label{subsec:ablation_design}

A central question left open by V2 is which components of \simon{} are responsible for orbital stability. We design a four-run ablation in which each run removes exactly one architectural component while keeping all others identical. The design is summarised in Table~\ref{tab:ablation}.

\begin{table}[H]
\centering
\caption{Four-run ablation framework. Each run removes exactly one component. All runs: IC1, masses $[1.0, 0.01, 0.005]\,M_\odot$, $T=100$\,yr, $dt=0.04$\,yr, same trained model weights. Final RMS vs \ias{} at $T=100$\,yr.}
\label{tab:ablation}
\begin{tabular}{llllcc}
\toprule
\textbf{Run} & \textbf{Adaptive stepping} & \textbf{Force direction} & \textbf{NN correction} & \textbf{Outcome} & \textbf{Final RMS} \\
\midrule
A & \textbf{OFF} & Analytic & ON & Ejection & 161.7\,AU \\
B & ON & \textbf{Learned (Vector NN)} & N/A & Ejection & 83.0\,AU \\
C & ON & Analytic & \textbf{OFF} ($c{=}1$) & Bounded & 1.67\,AU \\
D (\simon{}) & ON & Analytic & ON & Bounded & 1.54\,AU \\
\bottomrule
\end{tabular}
\end{table}

Runs A and B cause body ejection despite having adaptive stepping and NN correction (Run A) or adaptive stepping and analytic magnitude (Run B) active. Runs C and D maintain bounded dynamics. The contrast establishes that adaptive sub-stepping and analytic force direction are the stability-critical components. The $8.5\%$ difference between Runs C and D establishes the NN correction as an accuracy refinement rather than a stability requirement.

\subsubsection{Run A: Adaptive Sub-Stepping Is the Primary Stability Mechanism}
\label{subsec:run_a}

Run A tests the integrator with adaptive sub-stepping disabled. A frontier sweep across five timestep values ($dt \in [0.005, 0.01, 0.02, 0.04, 0.08]$\,yr) characterises behaviour across the tested range.

\begin{figure}[H]
    \centering
    \includegraphics[width=0.82\linewidth]{adaptive_error_vs_dt.png}
    \caption{Final RMS error vs timestep $dt$ before (red circles) and after (green squares) adding adaptive sub-stepping, $T=100$\,yr. Without adaptive sub-stepping, $dt=0.01$\,yr triggers body ejection to $161.7$\,AU -- a leapfrog failure where a close encounter is under-resolved within a single step. With adaptive sub-stepping, all five tested timestep values produce bounded dynamics. Worst-case improvement: $161.7$\,AU $\to 1.54$\,AU ($105\times$).}
    \label{fig:adaptive_frontier}
\end{figure}

Without adaptive sub-stepping, $dt=0.01$\,yr causes body ejection to $161.7$\,AU. With adaptive sub-stepping active, all five tested timestep values produce bounded simulations. The worst-case final RMS across the frontier reduces from $161.7$\,AU (ejection, before adaptive) to $1.54$\,AU (bounded, after adaptive): a $105\times$ improvement. The mechanism activates for fewer than $1\%$ of integration steps, targeting only close encounters below $0.05$\,AU.

At the operational timestep $dt=0.04$\,yr (the only point where \simon{} exceeds \ias{} in speed), adaptive sub-stepping reduces final RMS from $1.69$\,AU to $1.54$\,AU and divergence rate from $\lambda=0.1712$\,yr$^{-1}$ to $\lambda=0.1655$\,yr$^{-1}$, at a cost of $+28\%$ per-step overhead. Speedup vs \ias{} reduces from $1.53\times$ to $1.10\times$, remaining faster than the reference integrator.

\subsubsection{Run B: Analytic Force Direction Is Necessary for Stability}
\label{subsec:run_b}

Run B tests a Vector NN architecture in which both the force magnitude \emph{and} direction are predicted by the neural network, while adaptive sub-stepping remains active. The Vector NN achieves an average directional error of $0.04^\circ$ on the training distribution yet causes body ejection.

\begin{figure}[H]
    \centering
    \includegraphics[width=0.82\linewidth]{divergence_3model.png}
    \caption{RMS position error vs time comparing \simon{} (green solid, $\lambda=0.1655$\,yr$^{-1}$, bounded) and Vector NN (red dashed, diverging after $\sim60$\,yr to $>10^2$\,AU, ejection). Both models track closely for the first 60 years; after this, the Vector NN's accumulated non-radial impulses drive exponential divergence to body ejection.}
    \label{fig:run_b_divergence}
\end{figure}

\begin{figure}[H]
    \centering
    \includegraphics[width=0.95\linewidth]{trajectory_3model.png}
    \caption{XY trajectory overlay for all three bodies at $T=100$\,yr. \ias{} (blue solid) and \simon{} (green dashed) produce bounded overlapping orbits. Vector NN (red dotted) shows clear body ejection, with Body 2 reaching approximately $83$\,AU.}
    \label{fig:run_b_traj}
\end{figure}

The Vector NN final RMS of $83.0$\,AU represents body ejection despite adaptive stepping being active throughout. The root cause is the accumulation of small non-radial force errors at close encounters: over many orbital periods, these errors inject angular momentum into the system in an uncontrolled direction, destabilising the orbit. \simon{}'s analytic force direction -- using the exact normalised displacement $\mathbf{r}_{ij}/\|\mathbf{r}_{ij}\|$ -- prevents this accumulation entirely.

\subsubsection{Run C: The NN Correction Is an Accuracy Mechanism, Not a Stability Mechanism}
\label{subsec:run_c}

Run C tests the integrator with the NN correction disabled: the correction factor is fixed at $c = 1.0$ for all pair evaluations, reducing to pure softened gravity.

\begin{figure}[H]
    \centering
    \includegraphics[width=0.82\linewidth]{divergence_no_nn.png}
    \caption{RMS position error vs time for \simon{}/Run D (green solid, $\lambda=0.1655$\,yr$^{-1}$) and No-NN baseline/Run C (red dashed, $\lambda=0.1655$\,yr$^{-1}$). Both models produce nearly identical divergence curves throughout 100 years. The NN correction contributes an 8.5\% accuracy improvement in final RMS ($1.67$\,AU vs $1.54$\,AU) without affecting the divergence rate.}
    \label{fig:run_c_divergence}
\end{figure}

Run C maintains bounded dynamics (final RMS $= 1.67$\,AU) with an identical divergence rate ($\lambda = 0.1655$\,yr$^{-1}$, matching Run D to four decimal places). This directly answers the question left open by V2: the NN correction does not provide orbital stability. Adaptive sub-stepping alone (with analytic force direction) is sufficient for bounded long-horizon dynamics over 100 years. The NN contributes an $8.5\%$ improvement in final position accuracy at the $0.45\%$ of pair evaluations involving close encounters.

\subsubsection{Synthesis: Two Complementary Physical Constraints}
\label{subsec:synthesis}

The ablation results establish that \simon{}'s stability rests on two necessary and orthogonal constraints:

\begin{enumerate}
    \item \textbf{Adaptive sub-stepping} provides the primary stability guarantee. Without it (Run A), body ejection occurs at $dt=0.01$\,yr regardless of the NN correction.
    \item \textbf{Analytic force direction} prevents cumulative non-radial impulse injection. Without it (Run B), the Vector NN's small directional errors destabilise the orbit despite adaptive stepping being active.
\end{enumerate}

With both constraints active (Run D), \simon{} achieves bounded dynamics across all tested timestep values. The NN correction then provides an additional $8.5\%$ accuracy improvement. Figure~\ref{fig:simon_traj} shows the qualitative orbital structure for the full \simon{} system.

\begin{figure}[H]
    \centering
    \includegraphics[width=0.9\linewidth]{traj_overlay_xy_subplots_all_bodies_T100_after_adaptive.png}
    \caption{XY trajectory overlays for all three bodies over 100\,yr comparing \simon{} (dashed) to \ias{} (solid). \simon{} emulates the qualitative orbital geometry, including shape, scale, and boundedness of the \ias{} reference. The main discrepancy is phase drift accumulating after 25--30 years, consistent with the chaotic divergence inherent to the Three-Body problem.}
    \label{fig:simon_traj}
\end{figure}

\begin{figure}[H]
    \centering
    \includegraphics[width=0.9\linewidth]{traj_overlay_timeseries_all_bodies_T100_after_adaptive.png}
    \caption{Coordinate time series overlays $x(t)$, $y(t)$, for all three bodies over 100 years comparing \simon{} to \ias{}. For the first 25 years, body 0 phases follow \ias{} closely. After 25 years, amplitude and frequency continue to match but a phase offset develops. By year 60, curves are nearly half a period out of phase. The key observation is that all predicted \simon{} coordinates remain oscillatory and bounded within similar ranges to \ias{}. The primary error is accumulated phase drift \cite{ott2002chaos}, not instability.}
    \label{fig:simon_timeseries}
\end{figure}

\begin{figure}[H]
    \centering
    \includegraphics[width=0.82\linewidth]{model_vs_rebound_divergence_T100_after_adaptive.png}
    \caption{\simon{}--\ias{} RMS position deviation vs time over 100 years. The shaded window (10--50\,yr) is used for log-linear regression yielding $\lambda \approx 0.1655$\,yr$^{-1}$. The exponential growth with no discontinuities demonstrates that \simon{} reproduces the system's chaotic divergence structure without introducing catastrophic instabilities \cite{strogatz2015nonlinear,poincare1890sur}.}
    \label{fig:simon_divergence}
\end{figure}

% ---- 5.2 RESULTS II: GENERALISATION ----
\subsection{Results II: Generalisation Across Initial Conditions}
\label{sec:results_gen}

\subsubsection{Seven Initial Conditions Spanning the Dynamical Spectrum}

To assess generalisation, \simon{} is evaluated on seven initial conditions spanning the full range of three-body dynamical behaviour. Table~\ref{tab:multi_ic} summarises all results.

\begin{table}[H]
\centering
\caption{Multi-IC evaluation results. $dt=0.04$\,yr, $T=100$\,yr. $\lambda$ fitted over 10--50\,yr. Final RMS vs \ias{} at $T=100$\,yr. NN\_frac is fraction of pair evaluations where NN correction was applied.}
\label{tab:multi_ic}
\begin{tabular}{llccccc}
\toprule
\textbf{IC} & \textbf{Configuration} & $\boldsymbol{\lambda}$ \textbf{(yr$^{-1}$)} & \textbf{RMS (AU)} & \textbf{Bounded} & \textbf{Speedup} & \textbf{Note} \\
\midrule
IC1 & Default (V2 baseline)     & $0.1655$ & $1.543$  & YES & $1.11\times$ & \\
IC2 & Near-equal mass           & $0.0430$ & $116.87$ & NO  & $1.69\times$ & Unstable$^\dagger$ \\
IC3 & Tight inner pair          & $0.1207$ & $3.453$  & YES & $1.27\times$ & \\
IC4 & Hierarchical (outer 5 AU) & $0.0990$ & $1.598$  & YES & $1.12\times$ & \\
IC5 & High eccentricity         & $0.0648$ & $19.971$ & NO  & $1.17\times$ & Unstable$^\dagger$ \\
IC6 & Near-circular             & $-0.019$ & $0.008$  & YES & $1.07\times$ & Low chaos \\
IC7 & High ecc (redesign)       & $0.0341$ & $11.444$ & NO  & $1.07\times$ & Unstable \\
\midrule
\multicolumn{2}{l}{\textit{Bounded mean $\pm$ std (IC1,3,4,6)}} & $0.092\pm0.068$ & $1.65\pm1.22$ & 4/7 & & \\
\bottomrule
\end{tabular}
\\[0.5em]
{\small $^\dagger$\,Confirmed by independent \ias{} run: IC2 \ias{} final separation $58.9$\,AU; IC5 \ias{} ejects at $t=34.8$\,yr.}
\end{table}

\subsubsection{Bounded Configurations: Four Initial Conditions}

Four of the seven configurations produce bounded dynamics over 100 years. IC6 (near-circular) achieves near-zero divergence ($\lambda = -0.019$\,yr$^{-1}$) and $0.008$\,AU final RMS, demonstrating that \simon{} correctly preserves low-chaos configurations. IC3 (tight inner pair) shows higher divergence ($\lambda = 0.1207$\,yr$^{-1}$) due to frequent close encounters. IC4 (hierarchical) maintains bounded dynamics at $1.60$\,AU despite its IC4 energy drift discussed in Section~\ref{sec:energy}. All bounded ICs confirm that \simon{} generalises beyond the V2 default configuration.

\begin{figure}[H]
    \centering
    \includegraphics[width=0.9\linewidth]{multi_ic_divergence_v3.png}
    \caption{RMS position error vs time for all seven initial conditions. Bounded configurations (IC1, IC3, IC4, IC6) show chaotic saturation. Unstable configurations (IC2, IC5, IC7) show continued exponential divergence, indicating body ejection.}
    \label{fig:multi_ic_div}
\end{figure}

\begin{figure}[H]
    \centering
    \includegraphics[width=0.9\linewidth]{multi_ic_results_v3.png}
    \caption{Summary of final RMS (left) and divergence rate $\lambda$ (right) across all seven initial conditions. Bounded ICs (green) cluster below $4$\,AU final RMS. Unstable ICs (red) range from $11$\,AU to $117$\,AU.}
    \label{fig:multi_ic_results}
\end{figure}

\begin{figure}[H]
    \centering
    \includegraphics[width=0.82\linewidth]{traj_overlay2.pdf}
    \caption{Additional trajectory overlay for IC3 (tight inner pair). \simon{} displays consistent bounded performance over this setup, following a similar divergence-over-time trend to IC1.}
    \label{fig:traj2}
\end{figure}

\begin{figure}[H]
    \centering
    \includegraphics[width=0.82\linewidth]{traj_overlay3.pdf}
    \caption{Additional trajectory overlay for IC4 (hierarchical configuration). \simon{} displays consistent bounded performance, following a similar divergence-over-time trend. As we can see from these additional overlays, \simon{} is able to display consistent performance over various different setups, implying it is versatile and stable.}
    \label{fig:traj3}
\end{figure}

\subsubsection{Physically Unstable Configurations: Correct Reproduction of Instability}

Three configurations (IC2, IC5, IC7) produce final RMS values indicative of body ejection. For IC2 (near-equal mass, $[1.0, 0.5, 0.25]\,M_\odot$) and IC5 (high eccentricity), independent \ias{} runs confirm physical instability: \ias{} predicts ejection for both. For IC2, \simon{} predicts $116.9$\,AU versus \ias{}'s $58.9$\,AU -- a $2\times$ ratio consistent with chaotic amplification of integrator differences after the system becomes dynamically unbound. For IC5, \simon{} predicts $19.97$\,AU versus \ias{}'s ejection at $t=34.8$\,yr. IC7 (a redesigned high-eccentricity configuration) also produces an unbound trajectory ($11.44$\,AU), demonstrating that even with a more conservative periapsis, this parameter regime remains dynamically unstable.

Critically, \simon{} does not suppress physical instability. All three unstable configurations are correctly identified as unbound. The post-ejection separation differences reflect chaotic divergence amplifying integrator differences after dynamical unbinding -- not a failure of \simon{} to identify instability.

% ---- 5.3 RESULTS III: PHYSICAL CONSISTENCY ----
\subsection{Results III: Physical Consistency and Reproducibility}
\label{sec:results_phys}

\subsubsection{Energy Conservation: Leapfrog Attribution}
\label{sec:energy}

We compute the fractional energy drift $|\Delta E(t)/E_0| = |(E(t) - E(0))|/|E(0)|$ over 100 years at $dt=0.04$\,yr for three initial conditions (Table~\ref{tab:energy}).

\begin{table}[H]
\centering
\caption{Fractional energy drift at $T=100$\,yr for \simon{} and \ias{}. NN\_frac is the NN invocation rate. The proof: IC3 and IC4 have NN\_frac $= 0.000$ yet show energy drift comparable to IC1 -- confirming the drift originates from the leapfrog integrator, not the NN correction.}
\label{tab:energy}
\begin{tabular}{lccccc}
\toprule
\textbf{IC} & \textbf{$|E_0|$} & \textbf{\ias{} drift} & \textbf{\simon{} drift} & \textbf{NN\_frac} & \textbf{Source} \\
\midrule
IC1 (Default)     & $0.00723$ & ${\sim}10^{-14}\%$ & $1.39\%$  & $0.0045$ & Leapfrog + minor NN \\
IC3 (Tight pair)  & $0.01325$ & ${\sim}10^{-14}\%$ & $0.67\%$  & $0.0000$ & Leapfrog only \\
IC4 (Hierarchical)& $0.00602$ & ${\sim}10^{-14}\%$ & $13.77\%$ & $0.0000$ & Leapfrog (small $|E_0|$) \\
\bottomrule
\end{tabular}
\end{table}

\ias{} conserves energy to near float64 machine precision (${\sim}10^{-14}\%$). \simon{}'s energy drift is orders of magnitude larger, but its attribution is unambiguous: both IC3 and IC4 have NN\_frac $= 0.000$ -- the NN was never invoked. Both show drift comparable to IC1. The NN correction cannot cause energy drift it never contributed to. The drift originates from the second-order leapfrog integrator.

IC4's $13.77\%$ drift requires honest framing: it reflects the small binding energy $|E_0| = 0.006$ (the smallest of the three tested configurations), which amplifies the fractional representation of a leapfrog error that is comparable in absolute magnitude to IC1 and IC3. The drift is oscillatory rather than monotonic (consistent with symplectic shadow Hamiltonian conservation), and all bodies remain bounded (final RMS $= 1.60$\,AU). At $dt=0.01$\,yr, IC4's drift reduces to $1.05\%$, confirming the timestep dependence.

\begin{figure}[H]
    \centering
    \includegraphics[width=0.88\linewidth]{energy_drift_all_ics.png}
    \caption{Fractional energy drift $|\Delta E/E_0|$ over 100 years for \ias{} (near-flat, near machine precision) and \simon{} (coloured, characteristic leapfrog oscillation) across three initial conditions. IC3 and IC4 both have NN\_frac $= 0.000$; their comparable drift to IC1 confirms the leapfrog attribution.}
    \label{fig:energy_drift}
\end{figure}

\subsubsection{Timestep Scaling: Quantitative Confirmation of $dt^2$ Prediction}
\label{sec:dt2}

Second-order leapfrog theory predicts energy drift scaling as $dt^2$. Running \simon{} at $dt=0.01$\,yr and $dt=0.04$\,yr allows direct testing.

\begin{figure}[H]
    \centering
    \includegraphics[width=0.82\linewidth]{01_drift_scaling_loglog.png}
    \caption{Log-log plot of maximum fractional energy drift vs timestep for IC1, IC3, and IC4. Fitted slopes (IC1: $1.32$; IC3: $2.43$; IC4: $1.86$; average $1.87$) bracket the theoretical $dt^2$ prediction (grey reference line, slope $= 2.0$). The operational timestep $dt=0.04$\,yr is marked by the vertical dashed line. IC3 and IC4 (NN\_frac $= 0.000$) show the same scaling as IC1, providing independent quantitative confirmation that scaling is governed entirely by the leapfrog timestep.}
    \label{fig:dt_scaling}
\end{figure}

Fitted slopes of $1.32$, $2.43$, and $1.86$ (average $1.87$) are consistent with the theoretical prediction of $2.0$. IC3 and IC4 (NN\_frac $= 0.000$) follow the same scaling pattern as IC1, confirming quantitatively that the scaling is entirely determined by the leapfrog timestep. At $dt=0.01$\,yr, IC4's drift reduces from $13.77\%$ to $1.05\%$ (factor $13.2\times$; prediction: $16\times$).

\subsubsection{Computational Performance: Speed-Accuracy Frontier}
\label{sec:perf}

In this measurement, we observe the final position deviation of the bodies predicted by \simon{}, in comparison with \ias{}, over a 100 year time horizon. The x-axis is the throughput (how often the current state of the system is sampled) and the y-axis is RMS position error on a logarithmic scale. Points toward the lower-right are better (higher throughput, lower error). Throughput is computed as $n_{\text{samples}}/t_{\text{total}}$.

\begin{figure}[H]
    \centering
    \includegraphics[width=0.75\linewidth]{speed_accuracy_frontier_T100_after_adaptive.png}
    \caption{Speed--accuracy tradeoff across timesteps $dt$. \simon{} outperforms \ias{} at $dt=0.04$\,yr ($1.10\times$ speedup) and $dt=0.08$\,yr ($2.07\times$ speedup).}
    \label{fig:speed_acc}
\end{figure}

\begin{table}[H]
\centering
\caption{Speed-accuracy frontier results for \simon{} with adaptive sub-stepping, $T=100$\,yr. \ias{} baseline: $0.069$\,s. Bold row indicates the operational timestep.}
\label{tab:perf}
\begin{tabular}{ccccc}
\toprule
$\boldsymbol{dt}$ \textbf{(yr)} & \textbf{Final RMS (AU)} & \textbf{Throughput (samp/s)} & \textbf{Time/step ($\mu$s)} & \textbf{Speedup vs \ias{}} \\
\midrule
$0.005$ & $1.022$ & $10{,}333$  & $24.2$ & $0.14\times$ \\
$0.01$  & $0.847$ & $20{,}383$  & $24.5$ & $0.28\times$ \\
$0.02$  & $1.022$ & $40{,}119$  & $24.9$ & $0.55\times$ \\
$\mathbf{0.04}$  & $\mathbf{1.543}$ & $\mathbf{79{,}620}$  & $\mathbf{25.1}$ & $\mathbf{1.10{\times}}$ \\
$0.08$  & $1.022$ & $150{,}146$ & $26.6$ & $2.07\times$ \\
\bottomrule
\end{tabular}
\end{table}

The \ias{} baseline completes in $0.069$\,s. \simon{} outperforms \ias{} at $dt=0.04$\,yr ($1.10\times$ speedup) and $dt=0.08$\,yr ($2.07\times$ speedup). The lowest final RMS is $0.847$\,AU at $dt=0.01$\,yr, but this is $3.5\times$ slower than \ias{}. The errors at $dt=0.005$, $0.02$, and $0.08$ are all approximately $1.02$\,AU, reflecting chaotic saturation: after 100 years, any two different integrations of the same system will diverge to roughly one orbital radius regardless of timestep.

The per-step time ranges from $24.2\,\mu$s to $26.6\,\mu$s across a $16\times$ change in timestep -- a range of only $2.4\,\mu$s. This nearly flat profile confirms that runtime cost is dominated by per-step overhead (force evaluation, distances, adaptive sub-stepping) rather than integration arithmetic. Cost is dominated by force evaluation plus gating plus framework overhead, not integrator math. To improve speed, optimizing data movement is more important than tuning $dt$.

\begin{figure}[H]
    \centering
    \includegraphics[width=0.7\linewidth]{comp_cost_frontier_T100_after_adaptive.png}
    \caption{Speed-accuracy scatter plot: final RMS deviation vs throughput for each $dt$ value. The operational point $dt=0.04$\,yr achieves $1.10\times$ speedup over \ias{} with $1.54$\,AU final RMS.}
    \label{fig:comp_cost}
\end{figure}

\subsubsection{Training Seed Robustness}
\label{sec:seeds}

The PairCorrectionNN was trained five times with different random seeds ($[42, 123, 456, 789, 1000]$), all other hyperparameters held constant.

\begin{figure}[H]
    \centering
    \includegraphics[width=0.9\linewidth]{seed_robustness_results.png}
    \caption{Left: Validation loss curves over 5,000 epochs for all five seeds (log scale). All seeds converge smoothly without instability. Right: Best validation MSE per seed. Mean $= 1.54\times10^{-5}$, coefficient of variation $= 15.8\%$.}
    \label{fig:seed_robust}
\end{figure}

Validation MSE ranged from $1.11\times10^{-5}$ (seed 42, best) to $1.75\times10^{-5}$ (seed 456, worst), giving CV $= 15.8\%$. However, the divergence rate $\lambda = 0.1655$\,yr$^{-1}$ was identical to four decimal places across all five seeds, and all seeds produced bounded dynamics.

The 15.8\% CV does not imply dynamical sensitivity to seed choice: it is concentrated at sub-softening distances ($r < 5\times10^{-4}$\,AU) where the safety gate bypasses NN predictions entirely. In the operationally relevant range ($r \geq 10^{-3}$\,AU), all seeds produce correction factor errors below $0.3\%$, and downstream simulation metrics are identical across all seeds.

% ============================================================
% SECTION 6 -- DISCUSSION AND CONCLUSION (NEW)
% ============================================================
\section{Discussion and Conclusion}
\label{sec:discussion}

\subsection{The Two-Constraint Principle}

The central finding of this paper is that \simon{}'s long-horizon orbital stability requires two physically motivated constraints acting in concert. Adaptive sub-stepping resolves close-encounter instabilities in the fixed-timestep leapfrog scheme: without it, body ejection occurs regardless of the NN correction (Run A: $161.7$\,AU). Analytic force direction prevents the cumulative injection of non-radial angular momentum: without it, learned directional errors cause ejection despite adaptive stepping being active (Run B: $83.0$\,AU). Neither constraint alone is sufficient; together they guarantee bounded dynamics across all tested timestep values and initial conditions.

This decomposition has a broader implication for neural integrators in gravitational dynamics: the architecture must separately address the stability problem (which leapfrog alone cannot handle at close encounters) and the accuracy problem (which the NN addresses as a force magnitude correction). Conflating these roles -- as V2 implicitly did by framing the NN as handling ``rare edge cases that would otherwise cause failure'' -- produces a misleading picture of what is architecturally load-bearing.

\subsection{Positioning \simon{} Relative to \ias{}}

\simon{} does not compete with \ias{} on accuracy. \ias{} achieves near machine-precision energy conservation and is the gold standard for high-precision N-body integration. \simon{} targets a different operating point: long-horizon orbital characterisation where computational efficiency is prioritised. In a chaotic three-body system, exact trajectory reproduction becomes physically meaningless beyond the Lyapunov time ($1/\lambda \approx 6$\,yr for IC1): beyond this timescale, no integrator recovers the exact trajectory regardless of its order or precision. The relevant question is whether \simon{} correctly reproduces the qualitative dynamical behaviour (bounded vs unbounded, chaotic divergence rate, orbital geometry). Our results show it does, at $1.10\times$ the speed of \ias{} at the operational timestep.

\subsection{Limitations and Scope}

Several limitations should be noted. First, the ablation study (Section~\ref{sec:results_stab}) uses a single default initial condition (IC1). Generalisation is demonstrated across seven configurations in Section~\ref{sec:results_gen}, but the ablation claims are technically established for IC1 only. Second, the divergence-rate proxy $\lambda$ conflates intrinsic chaos with model mismatch and should not be interpreted as the formal Lyapunov exponent. Third, all mass ratios in the ablation and energy experiments use the V2 default configuration $[1.0, 0.01, 0.005]\,M_\odot$; performance for other mass regimes has not been systematically tested.

\subsection{Future Work}

Several natural extensions follow from this work. Evaluating \simon{} across broader mass ratios and initial condition distributions would establish more comprehensive generalisation claims. Extending to $N > 3$ bodies would require scaling the pairwise correction architecture. Comparison against equivariant neural networks that structurally enforce rotational symmetry would test whether force-direction constraints could be architecturally guaranteed rather than analytically imposed. Longer simulation horizons ($T > 100$\,yr) would test whether the stability demonstrated here persists over astrophysically relevant timescales.

\subsection{Research Question Answer}

\begin{quote}
\emph{How closely can physically constrained neural network force models approximate the long-term stability of traditional numerical integrators in three-body simulations, while operating at substantially higher computational speed?}
\end{quote}

We can say that a physically constrained neural network can approximate the qualitative characteristics of numerical integrators in three-body simulations. All bodies remain gravitationally bounded across four of seven tested initial conditions, with a $1.10\times$ speedup over \ias{} at the operational timestep. However, predictions will accumulate phase drift which leads to misalignment between trajectories, consistent with the chaotic sensitivity of the Three-Body Problem (see Section~\ref{Chaos}).

Stability is preserved through two necessary and complementary mechanisms: an analytic force direction that prevents non-radial impulse accumulation, and an adaptive sub-stepping mechanism that resolves close encounters that would otherwise cause leapfrog failure. The neural network force correction provides an additional $8.5\%$ accuracy improvement at close-encounter steps, but is not required for orbital stability. The dominant failure mode is accumulated phase drift rather than instability, and the positive $\lambda$ value confirms chaotic sensitivity consistent with the Three-Body problem.

On computational speed, the speed advantage arises from the hybrid architecture allowing the use of a simpler second-order integrator with a large, fixed timestep, resolved by adaptive sub-stepping during close encounters. Runtime cost is dominated by force evaluation and framework overhead rather than integration arithmetic.

The bottom line is that physically constrained neural force models like \simon{} can approximate long-term qualitative stability of the Three-Body Problem while offering meaningful speedups, provided both stability constraints -- adaptive sub-stepping and analytic force direction -- are applied correctly.

% ============================================================
% REFERENCES
% ============================================================
\bibliographystyle{plainnat}
\bibliography{V3_bib}

% ============================================================
% APPENDIX
% ============================================================
\appendix

\section{Initial Condition Coordinates}
\label{app:ics}

Table~\ref{tab:ic_coords} provides the full initial conditions for all seven configurations used in Section~\ref{sec:results_gen}.

\begin{table}[H]
\centering
\caption{Initial conditions for all seven configurations. Masses in $M_\odot$; positions in AU; velocities in AU/yr. $E_{\rm total}$ in simulation units ($G=1$).}
\label{tab:ic_coords}
\small
\begin{tabular}{lllcl}
\toprule
\textbf{IC} & \textbf{Masses ($M_\odot$)} & \textbf{Configuration} & \textbf{$E_{\rm total}$} & \textbf{Regime} \\
\midrule
IC1 (Default)      & $[1.0, 0.01, 0.005]$ & Bodies at $[0,0]$, $[1,0]$, $[-1,0]$    & $-0.00723$ & Moderate scattering \\
IC2 (Near-equal)   & $[1.0, 0.5, 0.25]$   & Bodies at $[0,0]$, $[1,0]$, $[-1,0]$    & $-0.73460$ & Strongly interacting \\
IC3 (Tight pair)   & $[1.0, 0.01, 0.005]$ & Bodies at $[0,0]$, $[0.5,0]$, $[-1,0]$  & $-0.01325$ & Tight inner binary \\
IC4 (Hierarchical) & $[1.0, 0.01, 0.005]$ & Bodies at $[0,0]$, $[1,0]$, $[5,0]$     & $-0.00602$ & Outer body at 5\,AU \\
IC5 (High ecc)     & $[1.0, 0.01, 0.005]$ & Body 1 at $[0.3,0]$, $v=2.2$ (85\% $v_{\rm esc}$) & $-0.01167$ & High eccentricity \\
IC6 (Circular)     & $[1.0, 0.01, 0.005]$ & Keplerian circular orbits                & $-0.00420$ & Near-circular \\
IC7 (High ecc v2)  & $[1.0, 0.01, 0.005]$ & Body 1 at $[0.6,0]$, $v=1.6$ ($1.24\times v_{\rm circ}$) & $-0.00520$ & High ecc (redesign) \\
\bottomrule
\end{tabular}
\end{table}

\section{Supplementary Figures}
\label{app:figs}

\begin{figure}[H]
    \centering
    \includegraphics[width=0.9\linewidth]{multi_ic_trajectories_v3.png}
    \caption{Trajectory overlays (\simon{} dashed, \ias{} solid) for all seven initial conditions. Bounded configurations (IC1, IC3, IC4, IC6) show qualitatively consistent bounded orbits. Unbounded configurations (IC2, IC5, IC7) show clear ejection trajectories.}
    \label{fig:app_traj}
\end{figure}

\begin{figure}[H]
    \centering
    \includegraphics[width=0.82\linewidth]{02_drift_timeseries_IC4_hierarchical.png}
    \caption{IC4 energy drift timeseries at five tested timestep values. At $dt=0.04$\,yr (operational, orange), drift reaches $13.77\%$ with oscillatory character (shadow Hamiltonian conserved). At $dt=0.01$\,yr (green), drift reduces to $1.05\%$. At $dt=0.08$\,yr and $dt=0.10$\,yr, integration breakdown occurs (permanent offset from $t\approx10$\,yr), confirming $dt=0.04$\,yr as the minimum viable operational timestep for IC4.}
    \label{fig:app_ic4}
\end{figure}

\end{document}
